\chapter{Cahaya sebagai Gelombang: Difraksi dan Interferensi}\label{ch:g12:light-as-wave}

Miringkanlah selaput sabun ke arah jendela: warna meluncur di
permukaannya, padahal tidak ada yang berwarna pada sabun atau airnya.
Arahkanlah \omterm{def:g11:color-light-sources:lamps}{laser} ke sehelai rambut: dinding di belakangnya menampilkan
bukan bayang-bayang setipis rambut melainkan sebuah goresan terang yang
tertutup garis gelap, ribuan kali lebih lebar daripada rambutnya. Sinar
yang berbaris lurus (\cref{ch:g10:refraction}) tidak menjelaskan kedua
penampakan itu. Bab ini menaikkan pangkat cahaya menjadi \omterm{def:g10:signals-and-waves:wave}{gelombang} yang
sepenuhnya --- setengah mikrometer dari puncak ke puncak --- lalu
menguangkan gagasannya: sebuah \omterm{def:g11:color-light-sources:lamps}{laser} dan sebuah dinding kosong akan
mengukur sehelai rambut, sebuah \omterm{def:g12:mechanical-waves:wavelength}{panjang gelombang}, dan alur sebuah CD.

\section{Cahaya adalah gelombang}

\begin{definition}[Gelombang cahaya]\label{def:g12:light-as-wave:lightwave}
\omterm{def:g10:light-spectra:white}{Cahaya monokromatik} --- satu warna murni, yang tidak dapat diuraikan
lagi oleh prisma mana pun (\cref{ch:g10:light-spectra}) --- adalah
\emph{gelombang cahaya}\index{gelombang cahaya} yang \omterm{def:g10:signals-and-waves:period}{periodik}: ia
menyeberangi ruang hampa dengan $c = \qty{3.00e8}{m/s}$ dan dicirikan
oleh \omterm{def:g10:signals-and-waves:frequency}{frekuensinya} $f$ atau oleh \emph{panjang gelombangnya di ruang
hampa}\index{panjang gelombang!cahaya} $\lambda = c/f$, yaitu jarak
puncak ke puncak pada \cref{ch:g12:mechanical-waves}. Mata menanggapi
dari sekitar \qty{400}{nm} (ungu) sampai \qty{800}{nm} (merah); \omterm{def:g10:light-spectra:white}{cahaya
putih} mencampurkan seluruh rentangnya.
\end{definition}

\begin{example}[Orde besar]\label{ex:g12:light-as-wave:orders}
\omterm{def:g11:color-light-sources:lamps}{Laser} helium--neon memancarkan cahaya merah dengan
$\lambda = \qty{633}{nm} =
\qty{6.33e-7}{m}$, sehingga
$f = c/\lambda \approx \qty{4.7e14}{Hz}$: tidak ada elektronik yang
sanggup mencacah secepat itu --- setiap \omterm{def:g12:mechanical-waves:wavelength}{panjang gelombang} pada bab ini
akan diukur dengan ilmu ukur. Sebagai pembanding, sehelai rambut
(sekitar \qty{70}{\micro m}) selebar seratus \omterm{def:g12:mechanical-waves:wavelength}{panjang gelombang}. Apa yang
bergelombang? Bukan udaranya: sebuah \omterm{def:g11:electric-gravitational-fields:efield}{medan listrik} dan \omterm{def:g11:magnetic-fields:bfield}{medan magnet},
yang dilukiskan secara jujur pada jilid Tahun ke-2; di bawah ini hanya
\omterm{def:g12:mechanical-waves:wavelength}{panjang gelombangnya} yang berperan.
\end{example}

\section{Difraksi}

\begin{definition}[Difraksi]\label{def:g12:light-as-wave:diffraction}
Ketika sebuah \omterm{def:g10:signals-and-waves:wave}{gelombang} menemui celah atau penghalang berukuran $a$
yang sebanding dengan \omterm{def:g12:mechanical-waves:wavelength}{panjang gelombangnya}, \omterm{def:g10:signals-and-waves:wave}{gelombang} itu menyebar ke
daerah yang akan ditinggalkan gelap oleh sinar. Penyebaran itu disebut
\emph{difraksi}\index{difraksi} --- yang dapat diabaikan selama
$\lambda/a$ sangat kecil, dan menguasai ketika $a$ menyusut menuju
$\lambda$.
\end{definition}

\begin{omfigure}
\begin{tikzpicture}[font=\small]
  \foreach \x in {-2.2,-1.75,...,-0.4}
    \draw[omDef, thick] (\x,-1.6) -- (\x,1.6);
  \draw[very thick] (0,-1.6) -- (0,-0.3) (0,0.3) -- (0,1.6);
  \foreach \r in {0.45,0.9,...,2.5}
    \draw[omDef, thick] (-75:\r) arc (-75:75:\r);
  \draw[Stealth-Stealth] (0.14,-0.3) -- (0.14,0.3)
    node[pos=0.28, right] {$a$};
\end{tikzpicture}

{\small \omterm{def:g12:mechanical-waves:wavefront}{Muka gelombang} datar (satu tiap puncak, terpisah $\lambda$)
menemui sebuah celah selebar $a$ yang sebanding dengan $\lambda$:
\omterm{def:g10:signals-and-waves:wave}{gelombangnya} terkipas keluar --- itulah \omterm{def:g12:light-as-wave:diffraction}{difraksi}.}
\end{omfigure}

\begin{proposition}[Setengah lebar sudut penyebarannya]\label{prop:g12:light-as-wave:halfwidth}
Sebuah celah selebar $a$, yang disinari \omterm{def:g12:mechanical-waves:wavelength}{panjang gelombang}
$\lambda < a$, menyebarkan cahayanya di sekitar arah datangnya di dalam
setengah lebar sudut $\theta$ (dalam radian); layar pada jarak $D$
menampilkan pita terang pusat selebar $L$, yaitu dua kali
$D \tan\theta \approx D\,\theta$ (untuk sudut yang kecil):
\[
\theta \approx \frac{\lambda}{a}, \qquad L = \frac{2\lambda D}{a}.
\]
\end{proposition}
\admitted

\begin{remark}[Bacalah rumusnya terbalik]\label{rem:g12:light-as-wave:honest}
Makin \emph{sempit} celahnya, makin \emph{lebar} penyebarannya ---
kebalikan ramalan sinarnya. Penurunan yang jujur (menjumlahkan \omterm{def:g10:signals-and-waves:wave}{gelombang}
kecil yang dipancarkan kembali di sepanjang celahnya) bersemayam pada
jilid Tahun ke-2; di sini rumusnya menjadi alat ukur.
\end{remark}

\begin{example}[Sebuah laser menembus celah]\label{ex:g12:light-as-wave:slit}
\omterm{def:g11:color-light-sources:lamps}{Laser} merah, $\lambda = \qty{633}{nm}$; celahnya
$a = \qty{0.10}{mm}$; layarnya pada $D = \qty{2.00}{m}$. Maka
$\theta = \num{6.33e-7} /
\num{1.0e-4} = \qty{6.3e-3}{rad}$ ---
sepertiga derajat, tetapi pita pusat di dindingnya selebar
$L = 2 \times \num{6.33e-7} \times 2.00 /
\num{1.0e-4} \approx \qty{2.5}{cm}$, yang diapit garis gelap dan pita samping yang lebih
samar. Separuhkanlah celahnya dan pitanya menjadi dua kali lipat.
\end{example}

\begin{omfigure}
\begin{tikzpicture}[font=\small]
  \fill[gray!55] (-3.2,-0.22) rectangle (-2.2,0.22);
  \node[above] at (-2.7,0.25) {laser};
  \draw[omThm, very thick] (-2.2,0) -- (0,0);
  \draw[very thick] (0,-1.5) -- (0,-0.1) (0,0.1) -- (0,1.5);
  \node[below] at (0,-1.5) {celah $a$};
  \draw[dashed] (0,0) -- (4.6,0);
  \draw[omThm] (0,0) -- (4.6,0.9) (0,0) -- (4.6,-0.9);
  \draw (1.5,0) arc (0:11:1.5) node[right, xshift=1pt] {$\theta$};
  \draw[very thick] (4.6,-1.85) -- (4.6,1.85);
  \draw[Stealth-Stealth] (4.9,-0.9) -- (4.9,0.9)
    node[midway, right] {$L$};
  \draw[Stealth-Stealth] (0,-2.1) -- (4.6,-2.1)
    node[midway, below] {$D$};
\end{tikzpicture}

{\small Percobaan celah tunggal: sebuah kerucut bersetengah sudut
$\theta
\approx \lambda/a$ melukiskan pita pusat selebar
$L = 2\lambda D/a$.}
\end{omfigure}

\begin{proposition}[Penghalang mendifraksikan seperti celah]\label{prop:g12:light-as-wave:hair}
Sebuah penghalang \omterm{def:g11:lenses-and-eye:lens}{tipis} yang tak tembus cahaya selebar $d$ --- kawat,
rambut --- menghasilkan, di luar berkas datangnya, pola yang sama
dengan celah selebar $d$ (asas Babinet, yang diturunkan pada jilid
Tahun ke-2).
\end{proposition}
\admitted

\begin{method}[Mengukur lebar yang sangat kecil]\label{met:g12:light-as-wave:hairwidth}
\begin{enumerate}
  \item Bentangkanlah rambutnya melintasi berkas \omterm{def:g11:color-light-sources:lamps}{lasernya}; pasanglah
        layar pada jarak $D$ yang terukur (beberapa \omterm{def:g10:orders-of-magnitude:unit}{meter}).
  \item Ukurlah lebar pita pusatnya $L$, dari garis gelap ke garis
        gelap.
  \item Karena $L = 2\lambda D/d$: $d = 2\lambda D/L$ --- makin halus
        rambutnya, makin besar polanya.
\end{enumerate}
\end{method}

\begin{remark}[Difraksi membatasi setiap alat]\label{rem:g12:light-as-wave:limits}
Cahaya memasuki kamera, teleskop atau mata melalui celah selebar $a$,
sehingga bahkan \omterm{def:g11:lenses-and-eye:lens}{lensa} yang sempurna pun mengaburkan sebuah titik menjadi
bercak berukuran sudut sekitar $\lambda/a$: rincian yang lebih rapat
melebur. Karena itu mikroskop cahaya tampak tidak memisahkan apa pun
yang jauh di bawah $\lambda \approx
\qty{0.5}{\micro m}$ (bakteri bisa,
virus tidak); cermin teleskop dibangun lebar karena alasan yang sama.
\end{remark}

\section{Interferensi}

\begin{definition}[Sumber koheren, beda lintasan]\label{def:g12:light-as-wave:coherent}
Dua sumber \omterm{def:g10:signals-and-waves:wave}{gelombang} disebut \emph{koheren}\index{sumber koheren}
apabila \omterm{def:g10:signals-and-waves:frequency}{frekuensinya} sama dan keduanya bergetar seirama --- pada
praktiknya, apabila keduanya merupakan dua salinan satu \omterm{def:g10:signals-and-waves:wave}{gelombang} yang
diperoleh dengan membelahnya. Untuk titik $M$ yang dicapai keduanya,
yang disebut \emph{beda \omterm{def:g10:relative-motion:trajectory}{lintasan}}%
\index{beda lintasan} adalah $\delta = S_2M - S_1M$.
\end{definition}

\begin{proposition}[Syarat interferensi]\label{prop:g12:light-as-wave:conditions}
Di tempat dua \omterm{def:g10:signals-and-waves:wave}{gelombang} \omterm{def:g12:light-as-wave:coherent}{koheren} berpanjang \omterm{def:g10:signals-and-waves:wave}{gelombang} $\lambda$
bertindihan, keduanya berjumlah --- keduanya
\emph{berinterferensi}\index{interferensi}:
\begin{itemize}
  \item $\delta = k\lambda$ ($k$ bulat): seirama, saling menguatkan ---
        interferensi \emph{konstruktif}\index{interferensi konstruktif}
        (terang);
  \item $\delta = (k + \tfrac12)\lambda$: puncak bertemu lembah, saling
        menghapus --- interferensi
        \emph{destruktif}\index{interferensi destruktif} (gelap).
\end{itemize}
\end{proposition}

\begin{proof}
Menggeser \omterm{def:g12:mechanical-waves:periodic}{gelombang periodik} sebanyak \omterm{def:g12:mechanical-waves:wavelength}{panjang gelombang} yang bulat tidak
mengubah apa pun: puncaknya bertumpuk pada puncak. Tambahan setengah
\omterm{def:g12:mechanical-waves:wavelength}{panjang gelombang} menempatkan puncak pada lembah: \omterm{def:g10:signals-and-waves:wave}{gelombang} yang sama
besar saling menghapus.
\end{proof}

\begin{example}[Percobaan dua celah Young]\label{ex:g12:light-as-wave:young}
Lubangilah dua celah halus $S_1$, $S_2$, yang terpisah sepersekian
milimeter, pada sebuah pelat yang tak tembus cahaya, lalu sinarilah
keduanya dengan satu \omterm{def:g11:color-light-sources:lamps}{laser}: jadilah dua \omterm{def:g12:light-as-wave:coherent}{sumber koheren}. Pada layar
beberapa \omterm{def:g10:orders-of-magnitude:unit}{meter} jauhnya, tindihannya bergaris-garis: di tengahnya
$\delta = 0$, terang; menaiki layarnya, $\delta$ tumbuh melalui
$\lambda/2$ (gelap), $\lambda$ (terang)\,\dots\ kegelapan di tempat
cahaya yang ditambahkan pada cahaya justru saling menghapus.
\end{example}

\begin{omfigure}
\begin{tikzpicture}[font=\small]
  \draw[very thick] (0,-1.7) -- (0,-0.72) (0,-0.48) -- (0,0.48)
    (0,0.72) -- (0,1.7);
  \fill (0,0.6) circle (1.5pt) node[left=2pt] {$S_1$} (0,-0.6) circle (1.5pt) node[left=2pt] {$S_2$};
  \draw[Stealth-Stealth] (-0.75,-0.6) -- (-0.75,0.6)
    node[midway, left] {$a$};
  \draw[very thick] (5.0,-2.0) -- (5.0,2.0);
  \draw[omDef] (0,0.6) -- (5.0,1.3) (0,-0.6) -- (5.0,1.3);
  \fill (5.0,1.3) circle (1.5pt) node[above right] {$M$};
  \draw[dashed] (0,0) -- (5.0,0) node[below right] {$O$};
  \draw[Stealth-Stealth] (5.55,0.04) -- (5.55,1.26) node[midway, right] {$x$};
  \draw[dashed] ([shift={(5.0,1.3)}]188:5.05) arc (188:200.8:5.05);
  \draw[omProp, very thick] (0,-0.6) -- (0.281,-0.493)
    node[below right=-2pt] {$\delta$};
  \foreach \y in {-1.8,-1.35,...,1.8}
    \fill[omExo] (5.06,\y-0.1) rectangle (5.2,\y+0.1);
\end{tikzpicture}

{\small Celah Young: \omterm{def:g10:relative-motion:trajectory}{lintasan} ke $M$ berbeda sebesar
$\delta = S_2M -
S_1M$; pita terangnya (goresan kelabu) berada di tempat
$\delta = k\lambda$.}
\end{omfigure}

\begin{definition}[Pita interferensi dan jarak antarpita]\label{def:g12:light-as-wave:fringes}
Pita terang dan gelap yang berselang-seling pada polanya disebut
\emph{pita interferensi}\index{pita interferensi}nya; jarak $i$ antara
pusat pita terang yang berurutan disebut \emph{jarak antarpita}%
\index{jarak antarpita}.
\end{definition}

\begin{proposition}[Jarak antarpita]\label{prop:g12:light-as-wave:interfringe}
Untuk celah yang terpisah sejauh $a$ dan layar pada jarak $D \gg a$,
titik pada layar yang berjarak $x$ dari pusat $O$ memiliki \omterm{def:g12:light-as-wave:coherent}{beda
lintasan} $\delta \approx a\,x/D$ (ilmu ukurnya diterima); menurut
\cref{prop:g12:light-as-wave:conditions}, pita terangnya berada di
$x_k = k\,\lambda D/a$, berjarak sama dengan \omterm{def:g12:light-as-wave:fringes}{jarak antarpita}
\[
i = \frac{\lambda D}{a}.
\]
Dengan $\lambda = \qty{633}{nm}$, $a = \qty{0.20}{mm}$,
$D =
\qty{2.0}{m}$: $i = \qty{6.3}{mm}$ --- susunannya memperbesar
$\lambda$ sebesar $D/a = 10^4$. Ilmu ukur di baliknya dikerjakan pada
jilid Tahun ke-2.
\end{proposition}
\admitted

\begin{omfigure}
\begin{tikzpicture}
\begin{axis}[omaxis, width=8.6cm, height=4.4cm,
    xlabel={$x$ (\unit{mm})}, ylabel={intensitas},
    xmin=-11, xmax=11, ymin=0, ymax=1.3, ytick=\empty]
  \addplot[omDef, thick, domain=-10.5:10.5, samples=211]
    {cos(deg(pi*x/3))^2 * exp(-x*x/50)};
  \draw[Stealth-Stealth] (axis cs:0,1.1) -- (axis cs:3,1.1)
    node[midway, above] {$i$};
\end{axis}
\end{tikzpicture}

{\small Intensitas melintasi layar pada percobaan Young: pitanya
terpisah $i = \lambda D/a$, di bawah \omterm{def:g10:refraction:fiber}{selubung} keseluruhan yang
digambarkan secara skematis.}
\end{omfigure}

\begin{method}[Mengukur sebuah panjang gelombang]\label{met:g12:light-as-wave:lambda}
\begin{enumerate}
  \item Kirimkanlah \omterm{def:g11:color-light-sources:lamps}{laser} yang belum diketahui itu menembus celah ganda
        yang jarak celahnya $a$ diketahui; tempatkanlah layarnya pada
        jarak $D$ yang terukur.
  \item Ukurlah rentang $10$ \omterm{def:g12:light-as-wave:fringes}{jarak antarpita}; bagilah dengan $10$:
        itulah \omterm{def:g12:light-as-wave:fringes}{jarak antarpitanya} $i$, dengan galat penggarisnya
        menyusut sepuluh kali.
  \item Lalu $\lambda = i\,a/D$.
\end{enumerate}
\end{method}

\begin{remark}[Mengapa dua lampu tidak pernah berinterferensi]\label{rem:g12:light-as-wave:lamps}
Dua lampu meja yang menyinari satu dinding tidak menghasilkan pita:
masing-masing memancarkan rangkaian \omterm{def:g10:signals-and-waves:wave}{gelombang} yang pendek dan tidak
berkorelasi, yang penjadwalannya melompat secara acak semiliar kali tiap
sekon, sehingga pitanya bergeser secepat itu pula dan mata
merata-ratakannya menjadi cahaya yang seragam. Interferensi menuntut
kekoherenan --- satu \omterm{def:g10:signals-and-waves:wave}{gelombang} yang dibelah menjadi dua, sebagaimana
dilakukan celah Young; \omterm{def:g11:color-light-sources:lamps}{laser} adalah jalan pintas modernnya, yaitu satu
\omterm{def:g10:signals-and-waves:wave}{gelombang} bersih yang panjang.
\end{remark}

\section{Warna yang dilukis interferensi}

\begin{definition}[Iridesensi]\label{def:g12:light-as-wave:iridescence}
Warna yang bergeser bersama sudut pandangnya --- gelembung sabun, lapisan
minyak, bulu merak, punggung sebuah CD --- disebut \emph{iridesensi}%
\index{iridesensi}: tanpa zat warna, melainkan interferensi di dalam
\omterm{def:g10:light-spectra:white}{cahaya putih}, yang menguatkan sebagian \omterm{def:g12:mechanical-waves:wavelength}{panjang gelombangnya} dan
menghapus yang lain.
\end{definition}

\begin{example}[Selaput sabun dan lapisan minyak]\label{ex:g12:light-as-wave:soap}
Cahaya yang mengenai selaput sabun sebagian memantul dari muka depannya,
sebagian dari muka belakangnya: dua salinan \omterm{def:g12:light-as-wave:coherent}{koheren} yang \omterm{def:g12:light-as-wave:coherent}{beda
lintasannya} ditetapkan tebalnya dan sudut pandangnya. Tiap tebalnya
mengembalikan sebagian \omterm{def:g12:mechanical-waves:wavelength}{panjang gelombang} dalam keadaan terkuatkan dan
sebagian lagi terhapus: selaputnya memantulkan warna yang meluncur
selagi ia meniris dan menipis. Lapisan minyak di atas aspal yang basah
memainkan muslihat yang sama.
\end{example}

\begin{remark}[Monokromatik lawan cahaya putih]\label{rem:g12:light-as-wave:white}
Di dalam \omterm{def:g10:light-spectra:white}{cahaya putih}, tiap \omterm{def:g12:mechanical-waves:wavelength}{panjang gelombang} melukis pita Young-nya
sendiri dengan $i = \lambda D/a$ miliknya sendiri: di tengahnya semua
warnanya terang (satu pita putih); beberapa pita ke luar, polanya
saling menjauh --- tepi yang berwarna pelangi, lalu kabur. \omterm{def:g10:light-spectra:white}{Cahaya
monokromatik} memberikan pita yang tajam tanpa henti: dari sinilah \omterm{def:g11:color-light-sources:lamps}{laser}
pada setiap pengukuran yang tajam di sini.
\end{remark}

\section{Latihan}

\begin{exercise}[$\star$]\label{exo:g12:light-as-wave:1}
Hitunglah \omterm{def:g10:signals-and-waves:frequency}{frekuensi} cahaya merah ($\lambda = \qty{700}{nm}$) dan cahaya
ungu (\qty{400}{nm}). Yang mana yang lebih besar?
\end{exercise}

\begin{exercise}[$\star$]\label{exo:g12:light-as-wave:2}
Sebuah \omterm{def:g11:color-light-sources:lamps}{laser} hijau ($\lambda = \qty{532}{nm}$) melewati celah selebar
$a = \qty{0.20}{mm}$: hitunglah $\theta$, lalu lebar pita pusatnya pada
layar yang berjarak \qty{1.5}{m}.
\end{exercise}

\begin{exercise}[$\star$]\label{exo:g12:light-as-wave:3}
Sebuah pintu selebar \qty{0.80}{m}: hitunglah $\lambda/a$ untuk suara
\qty{425}{Hz} ($v = \qty{340}{m/s}$) dan untuk cahaya \qty{550}{nm}.
Mengapa kamu dapat mendengar, tetapi tidak dapat melihat, mengitari
sudut ruangan?
\end{exercise}

\begin{exercise}[$\star$]\label{exo:g12:light-as-wave:4}
Dua pengeras suara yang digerakkan satu pembangkit memancarkan, secara
seirama, \omterm{def:g10:signals-and-waves:sound}{bunyi} berpanjang \omterm{def:g10:signals-and-waves:wave}{gelombang} \qty{2.0}{cm}. \omterm{def:g12:light-as-wave:coherent}{Beda lintasannya} di
$M$: \qty{6.0}{cm}; di $N$: \qty{5.0}{cm}. Nyaring atau senyap pada
masing-masing? Berikan alasannya.
\end{exercise}

\begin{exercise}[$\star$]\label{exo:g12:light-as-wave:5}
Celah Young: $\lambda = \qty{633}{nm}$, $a = \qty{0.25}{mm}$,
$D = \qty{1.8}{m}$. Hitunglah \omterm{def:g12:light-as-wave:fringes}{jarak antarpitanya}. Berapa pita terang
yang muat di dalam \qty{2.0}{cm} bagian tengah layarnya?
\end{exercise}

\begin{exercise}[$\star\star$]\label{exo:g12:light-as-wave:6}
Sehelai rambut di dalam berkas $\lambda = \qty{650}{nm}$, dengan layar
pada \qty{1.6}{m}, memberikan pita pusat selebar \qty{2.8}{cm}.
Hitunglah diameter rambutnya.
\end{exercise}

\begin{exercise}[$\star\star$]\label{exo:g12:light-as-wave:7}
Sebuah \omterm{def:g11:color-light-sources:lamps}{laser} yang belum diketahui melewati celah Young
($a = \qty{0.20}{mm}$, $D = \qty{1.50}{m}$): sepuluh \omterm{def:g12:light-as-wave:fringes}{jarak antarpita}
membentang \qty{4.7}{cm}. Tentukan $\lambda$ dan warnanya. Mengapa
mengukur sepuluh dan bukan satu?
\end{exercise}

\begin{exercise}[$\star\star$]\label{exo:g12:light-as-wave:8}
Dua lampu meja yang serupa menyinari dinding yang sama; berkasnya
bertindihan, namun tidak pernah ada pita yang muncul. Syarat kekoherenan
yang mana (\cref{def:g12:light-as-wave:coherent}) yang gagal, dan
bagaimana susunan Young memperbaikinya?
\end{exercise}

\begin{exercise}[$\star\star$]\label{exo:g12:light-as-wave:9}
Pada pola dua celah ($a = \qty{0.30}{mm}$, $D = \qty{1.4}{m}$), pusat
pita terang pertama dan ketujuh terpisah \qty{18.0}{mm}: tentukan \omterm{def:g12:light-as-wave:fringes}{jarak
antarpitanya}, lalu $\lambda$.
\end{exercise}

\begin{exercise}[$\star\star$]\label{exo:g12:light-as-wave:10}
Alur sebuah CD membentuk kisi berjarak $d = \qty{1.6}{\micro m}$. Dengan
menerima $d \sin\theta = k\lambda$ ($k$ bulat), tentukan arah berkasnya
untuk $\lambda = \qty{633}{nm}$. Berapa orde $k$ yang tertinggi?
\end{exercise}

\begin{exercise}[$\star\star$]\label{exo:g12:light-as-wave:11}
Sebuah teleskop bercelah $a = \qty{10}{cm}$ pada
$\lambda =
\qty{550}{nm}$: hitunglah kaburan \omterm{def:g12:light-as-wave:diffraction}{difraksinya}
$\theta \approx
\lambda/a$. Dapatkah teleskop itu memisahkan dua lampu
depan yang terpisah \qty{1.5}{m} pada jarak \qty{100}{km}?
\end{exercise}

\begin{exercise}[$\star\star\star$]\label{exo:g12:light-as-wave:12}
Sebuah selaput sabun yang tegak meniris, menjadi lebih tebal di bagian
bawahnya, lalu menampilkan pita warna mendatar yang merayap turun.
Jelaskan: mengapa ada warna, mengapa berupa pita, dan mengapa bergerak.
(Tanpa perhitungan.)
\end{exercise}

\begin{exercise}[$\star\star\star$]\label{exo:g12:light-as-wave:13}
Celah Young di dalam \omterm{def:g10:light-spectra:white}{cahaya putih}, $a = \qty{0.20}{mm}$,
$D = \qty{1.5}{m}$: hitunglah \omterm{def:g12:light-as-wave:fringes}{jarak antarpita} untuk ungu
(\qty{400}{nm}) dan merah (\qty{750}{nm}), lalu simpulkan apa yang
ditampilkan layarnya di tengahnya, beberapa milimeter ke luar, dan jauh
ke luar.
\end{exercise}

\begin{exercise}[$\star\star\star$]\label{exo:g12:light-as-wave:14}
Sebuah \omterm{def:g10:universal-gravitation:satellite}{satelit} pengintai membawa cermin \qty{2.4}{m} pada ketinggian
\qty{250}{km}. Dengan memakai $\theta \approx \lambda/a$ pada
\qty{550}{nm}, tentukan rincian terkecil di tanah yang dapat
dipisahkannya; nilailah pengakuan bahwa \omterm{def:g10:universal-gravitation:satellite}{satelit} itu ``membaca surat
kabar dari balik bahumu''.
\end{exercise}

\begin{exercise}[$\star\star\star$]\label{exo:g12:light-as-wave:15}
Sebuah susunan Young menampilkan $i = \qty{4.0}{mm}$ di udara, lalu
dibenamkan ke dalam air ($n = 1.33$). Dengan mengingat
\cref{ch:g10:refraction}: apakah \omterm{def:g10:signals-and-waves:frequency}{frekuensinya} berubah? Apakah \omterm{def:g12:mechanical-waves:wavelength}{panjang
gelombangnya} berubah? Hitunglah \omterm{def:g12:light-as-wave:fringes}{jarak antarpitanya} yang baru.
\end{exercise}

\section{Soal: Mengukur dengan cahaya}

\begin{problem}[{Soal akhir pekan --- mengukur dengan cahaya: sebuah
laser, sebuah penggaris dan sebuah dinding kosong menjadi bengkel
mikrometer yang menakar sebuah celah, sehelai rambut manusia, sebuah
panjang gelombang yang belum diketahui, dan alur dua cakram perak}]
\label{pb:g12:light-as-wave:1}
Perlengkapanmu: sebuah \omterm{def:g11:color-light-sources:lamps}{laser} merah ($\lambda = \qty{633}{nm}$, menurut
labelnya), sebuah \omterm{def:g11:color-light-sources:lamps}{laser} hijau yang \omterm{def:g12:mechanical-waves:wavelength}{panjang gelombangnya} belum
diketahui, sebuah celah tertera $a = \qty{100}{\micro m}$, sebuah celah
ganda berjarak $a' =
\qty{0.250}{mm}$, sebuah meteran gulung, dua helai
rambut, sebuah CD dan sebuah DVD.

\medskip
\noindent\textbf{Bagian I --- Menera pada celah yang diketahui.}
\omterm{def:g11:color-light-sources:lamps}{Laser} merah, celahnya, dan dinding pada $D = \qty{3.00}{m}$.
\begin{enumerate}
  \item Hitunglah setengah lebar sudut $\theta$ berkas yang
        terdifraksi.
  \item Tunjukkan bahwa pita pusatnya selebar $L = 2\lambda D/a$; lalu
        hitunglah nilainya.
  \item Kamu mengukur $L = \qty{3.7}{cm}$: hitunglah simpangan
        relatifnya. Apakah caranya sudah terbukti?
  \item Sebuah celah kedua lebarnya separuhnya, \qty{50}{\micro m}:
        ramalkan $L$-nya. Nyatakan aturan yang mengaitkan lebar celahnya
        dengan lebar polanya.
  \item Benarkan langkah $\tan\theta \approx \theta$ pada pertanyaan 2:
        bandingkan $\theta$ dan $\tan\theta$ pada sudut ini.
\end{enumerate}

\medskip
\noindent\textbf{Bagian II --- Rambutnya.} \omterm{def:g11:color-light-sources:lamps}{Laser} dan dinding yang sama;
sehelai rambut yang dibentangkan menggantikan celahnya.
\begin{enumerate}[resume]
  \item Hasil terima yang mana yang memungkinkanmu tetap memakai rumus
        celahnya untuk sehelai rambut, dan apa \omterm{def:g10:signals-and-waves:sound}{bunyinya}?
  \item Pita pusatnya selebar $L = \qty{5.4}{cm}$: hitunglah diameter
        rambutnya $d$.
  \item Penggarismu membaca $L$ sampai $\pm\qty{0.2}{cm}$: berikan
        rentang yang dihasilkannya untuk $d$, sebagai $d \pm \Delta d$.
  \item Rambut temanmu memberikan $L = \qty{7.6}{cm}$: berapa
        diameternya? Rambut siapa yang lebih halus, dan mengapa lebih
        halus berarti lebih lebar?
  \item Dapatkah sebuah lampu meja menggantikan \omterm{def:g11:color-light-sources:lamps}{lasernya}? Berikan dua
        alasan.
\end{enumerate}

\medskip
\noindent\textbf{Bagian III --- \omterm{def:g12:mechanical-waves:wavelength}{Panjang gelombang} \omterm{def:g11:color-light-sources:lamps}{laser} hijaunya.}
\omterm{def:g11:color-light-sources:lamps}{Laser} hijau, celah ganda, dan dinding pada $D = \qty{2.00}{m}$.
\begin{enumerate}[resume]
  \item Mengapa kedua celahnya harus disinari \omterm{def:g11:color-light-sources:lamps}{laser} yang \emph{sama}?
  \item Mengapa pita pusatnya terang berapa pun \omterm{def:g12:mechanical-waves:wavelength}{panjang gelombangnya}?
  \item Sepuluh \omterm{def:g12:light-as-wave:fringes}{jarak antarpita} membentang \qty{4.26}{cm}. Berikan
        $i$-nya, beserta alasan mengukur sepuluh sekaligus.
  \item Simpulkan \omterm{def:g12:mechanical-waves:wavelength}{panjang gelombang} \omterm{def:g11:color-light-sources:lamps}{laser} hijaunya.
  \item Label pembuatnya menyatakan $\qty{532 \pm 10}{nm}$: apakah
        sesuai? Warna apa yang kamu harapkan?
\end{enumerate}

\medskip
\noindent\textbf{Bagian IV --- Membaca cakram perak.} \omterm{def:g11:color-light-sources:lamps}{Laser} merahnya
mengenai sebuah CD tegak lurus permukaannya; alurnya berlaku sebagai
kisi pantul yang menaati (diterima) $d \sin\theta = k\lambda$. Pada
dinding di $D =
\qty{1.00}{m}$, tiap berkas orde pertamanya membuat
bercak pada jarak $x$ dari bercak pusatnya.
\begin{enumerate}[resume]
  \item Untuk CD-nya, $x = \qty{43}{cm}$: hitunglah $\theta$. Mengapa
        jalan pintas sudut kecil pada Bagian I kini terlarang?
  \item Simpulkan jarak alur CD-nya $d$.
  \item DVD-nya mengirimkan bercak orde pertamanya ke
        $x = \qty{1.65}{m}$: hitunglah $\theta$-nya dan jarak alurnya.
  \item Bandingkan kedua jarak alurnya. Mengapa DVD menyimpan lebih
        banyak daripada CD?
  \item Tutuplah buku catatanmu: daftarkan panjang yang terukur akhir
        pekan ini, lalu nyatakan dalam satu kalimat apa yang menjadikan
        cahaya sebagai penggaris bagi semuanya.
\end{enumerate}
\end{problem}
